% 摘要与讲述主线；技术规划完成审阅前保留槽位。
\section{项目摘要}
\KeyPoint{\StoryNineTitle}{\StoryNineQuestion}

\subsection{项目定位}
\noindent\begin{tabularx}{\linewidth}{@{}lX@{}}
\toprule
项目名称 & \ProjectName \\
企业命题 & \ChallengeTitle \\
命题企业 & \ChallengeCompany \\
参赛类别 & \CompetitionTrack{}／\CompetitionGroup \\
\bottomrule
\end{tabularx}

\ContentSlot{一页摘要}{本段应连贯说明企业要解决的问题、拟形成的系统、主要工作边界、验证方式及预期应用价值。}

\subsection{完成状态与阅读主线}
\KeyPoint{\StoryTenTitle}{\StoryTenQuestion}
\ContentSlot{现状与目标}{分别列出已有资源、设计中内容、拟实现功能和待验证结果，说明后续章节之间的关系。这里应与配置项和正文的完成状态保持一致。}
